\section{Introduction}
\label{sec:introduction}

If a linear map on a finite-dimensional vector space has a left inverse, that inverse is also a right inverse. A group algebra need not have finite dimension, but each of its elements still has finite support. Kaplansky's direct-finiteness conjecture asks whether this finite support is enough: for every field $K$, group $G$, and $a,b\in K[G]$, must $ab=1$ imply $ba=1$? Here $K[G]$ consists of finite formal sums $\sum_g c_g g$, multiplied using the group law.

A unital ring with this property is \emph{directly finite}; it is \emph{stably finite} if all its finite square matrix rings are directly finite. This distinction matters to a counterexample. Producing matrices $AB=I\ne BA$ over a group algebra does not by itself produce scalar elements with the same property. Our construction first produces the matrix pair, then embeds its matrix algebra in a larger group algebra to obtain the following scalar conclusion.

\begin{theorem}
\label{thm:main}
There exist a finite field $K$ of characteristic two, a finitely presented group $G$ containing an element of odd prime order, and elements $a_{\mathrm{out}},b_{\mathrm{out}}\in K[G]$ such that
\[
a_{\mathrm{out}}b_{\mathrm{out}}=1,
\qquad b_{\mathrm{out}}a_{\mathrm{out}}\ne1.
\]
The field, group, and finite sums are specified by a terminating prescription involving finite sets and finite-field arithmetic.
\end{theorem}

Thus Theorem~\ref{thm:main} resolves the unrestricted conjecture negatively in characteristic two. The construction uses odd-order torsion and does not assert that $K=\F_2$. The prescription need not be computationally practical and has not been executed: its termination is proved, and the two product assertions have algebraic certificates independent of any decision procedure for the group word problem.

\paragraph{History and positive results.}
For a finite group $G$, the algebra $K[G]$ and all its finite square matrix rings are finite dimensional, so elementary linear algebra gives stable finiteness. The classical characteristic-zero theorem uses a different substitute for dimension. Kaplansky's operator-algebra argument places $\mathbb C[G]$ in a finite tracial setting; Montgomery gave a proof using the operator-norm closure of the regular representation and its identity-coefficient trace, including the extension to every finite matrix size~\cite{Montgomery1969}. The trace rules out a proper one-sided inverse. For an arbitrary characteristic-zero field, all coefficients of a proposed finite matrix pair lie in a finitely generated subfield, which embeds in $\mathbb C$. Thus the complex result gives stable finiteness over every characteristic-zero field.

In positive characteristic, rank and approximation methods have supplied broad classes of positive results without a complex representation. Ara, O'Meara, and Perera proved stable finiteness over arbitrary division rings for free-by-amenable groups, using Sylvester rank functions and amenable translation rings~\cite{AraOMearaPerera2002}. Elek and Szab\'o obtained stable finiteness over arbitrary division rings for all sofic groups from finite approximate models and a rank construction~\cite[Corollary~4.7]{ElekSzabo2004}. These are restrictions on a possible counterexample's group, not assumptions that the group constructed here satisfies.

Dykema and Juschenko clarified how a matrix obstruction can nevertheless yield a scalar one: for a fixed field and group, stable finiteness is equivalent to direct finiteness after taking the direct product with every finite group~\cite[Theorem~2.2]{DykemaJuschenko2015}. Their reduction embeds a matrix algebra in a finite-group algebra and adds the complementary identity. We realize this transfer explicitly for the field used here, giving both the matrix units and the resulting two elements.

\paragraph{Finite presentation and soficity.}
For a finite nonempty set $\Omega$, write $d_\Omega(f,h)$ for the fraction of points on which two maps $f,h:\Omega\to\Omega$ differ. In the finite-map formulation of Elek and Szab\'o~\cite[Definition~4.2]{ElekSzabo2004}, a group $\Gamma$ is \emph{sofic} if, for every finite $E\subseteq\Gamma$ and every $\varepsilon\in(0,1)$, there are a finite nonempty $\Omega$ and a map $\sigma:\Gamma\to\Omega^\Omega$ such that
\[
\begin{aligned}
d_\Omega\bigl(\sigma(gh),\sigma(g)\sigma(h)\bigr)&\le\varepsilon
&& (g,h\in E),\\
d_\Omega\bigl(\sigma(1),\id\bigr)&\le\varepsilon,\\
d_\Omega\bigl(\sigma(g),\id\bigr)&>1-\varepsilon
&& (g\in E\setminus\{1\}),
\end{aligned}
\]
where products of maps denote composition. Thus soficity asks for finite approximate models of multiplication in which nonidentity elements act almost without fixed points.

The stable-finiteness theorem of Elek and Szab\'o~\cite[Corollary~4.7]{ElekSzabo2004} applies to the field $K$ in Theorem~\ref{thm:main}. If its particular group $G$ were sofic, then $K[G]$ would be stably finite and hence directly finite by the matrix-size-one case, contradicting the two displayed product assertions. Therefore this same $G$ is nonsofic. In particular, the construction supplies a finitely presented counterexample to the assertion that every group is sofic. The finite presentation is established directly in Section~\ref{sec:finite-presentation}; the nonsoficity conclusion uses only the scalar defect and the cited theorem.

\paragraph{Cellular automata.}
For a nonempty finite alphabet $A$, a configuration is an arbitrary function $x:G\to A$. A \emph{cellular automaton} on $A^G$ is a map given by a finite set $M\subset G$ and a local rule $\phi:A^M\to A$ through
\[
  \tau(x)(g)=\phi\bigl((x(gu))_{u\in M}\bigr).
\]
The group $G$ is \emph{surjunctive} if every injective cellular automaton on $A^G$ is surjective, for every such $A$. Gottschalk's conjecture asserts that every group is surjunctive~\cite{Gottschalk1973}. The problem asks whether finite local dependence and translation symmetry preserve the finite-set implication from injectivity to surjectivity on the infinite configuration space.

Gromov developed a general symbolic-algebraic surjunctivity theorem~\cite[Section~7.G$'$]{Gromov1999}; Weiss gave the finite-alphabet sofic-group theorem directly~\cite[Theorem~3.2]{Weiss2000}. Elek and Szab\'o explicitly recorded the implication from surjunctivity to direct finiteness over finite fields~\cite[Section~1]{ElekSzabo2004}. Ceccherini-Silberstein and Coornaert developed the linear theory: their equivalence, quantified over all coefficient fields, relates stable finiteness to injective-to-surjective behavior of linear cellular automata with finite-dimensional vector alphabets~\cite[Theorem~4.2]{CSCCoornaert2006}. Phung proved stable finiteness over every field for surjunctive groups, using algebraic cellular automata~\cite[Theorem~B]{Phung2023}. Ceccherini-Silberstein, Coornaert, and Phung later gave a model-theoretic proof and a monoid extension~\cite[Theorem~1.1]{CSCPhungGroup2025,CSCPhungMonoid2025}. These results connect algebraic finiteness to a dynamical property rather than imposing finite dimensionality.

For the finite field in Theorem~\ref{thm:main}, the short algebra-to-automaton passage can be given directly. The coefficients of $b_{\mathrm{out}}$ define a local rule, the coefficients of $a_{\mathrm{out}}$ give a left inverse on all configurations, and a nonzero coefficient of $b_{\mathrm{out}}a_{\mathrm{out}}-1$ identifies a missing point-indicator configuration. Section~\ref{sec:cellular} proves this standard transfer in full, including the multiplication convention. It gives a negative resolution of Gottschalk's conjecture for the group in Theorem~\ref{thm:main}; it requires no further existence theorem or approximation assumption.

The characteristic-two scalar pair in Theorem~\ref{thm:main} is also the
input to the companion counterexample to the group-ring Determinant
Conjecture~\cite[Theorem~1.1]{DeterminantCompanion}. By a separate
construction, that paper produces an integral square matrix over another
finitely generated group, invertible over its rational group ring, whose
Fuglede--Kadison determinant lies strictly between zero and one.

\paragraph{The construction.}
The obstacle is to force a one-sided inverse while retaining a certificate that the reverse product is not the identity. The proof separates these requirements through finite incidence data. Let $\mathcal F$ be a family of subsets of a finite set $V$. We require that each point lie in $\ell+1$ members, where $\ell$ is an odd prime, and that the bipartite incidence graph have girth at least $12$. We also require column vectors $x_v\in K^m$ satisfying
\[
\sum_{v\in f}x_vx_v^{\mathsf T}=0\quad(f\in\mathcal F),
\qquad \sum_{v\in V}x_vx_v^{\mathsf T}=-I_m,
\qquad m<|V|\le m\ell^2.
\]
Theorem~\ref{thm:criterion} turns these finite data into a counterexample.

Two parts of this criterion explain the separation. First, a planar balance argument proves that copies of the additive group $\F_\ell^2$, glued along their one-dimensional subgroups according to the incidence structure, embed in the resulting group. Incidence girth forces a contradiction with Euler's formula for any diagram witnessing a collapsed local element. Second, character idempotents and prescribed subgroup conjugations pack a rectangular factorization into square matrices. Their forward product is the identity by the displayed pairings.

To detect the reverse-product defect, we sandwich it by the packing maps and obtain a matrix in an embedded subgroup algebra. A homomorphism from that algebra to $K$ sends this matrix to $X^{\mathsf T}X-I_{|V|}$, where $X$ has columns $x_v$. Since $X$ has only $m<|V|$ rows, $X^{\mathsf T}X$ cannot be the identity, so the original defect is nonzero. The homomorphism is used only on the subgroup algebra; it is not asserted to extend to the larger group algebra. A finite auxiliary group then converts the matrix obstruction into the scalar outputs of Theorem~\ref{thm:main}.

To produce the incidence data, we take full affine lines over a finite field $Q\subset K$ and remove a small set of points. Low-degree interpolation on these holes gives the vectors $x_v$: the pairings sum to zero on each full line, but removing the holes leaves the nonzero total $-I_m$. The geometric task is to retain a regular line system of large incidence girth. We combine random small-dimensional slabs with an exact trade in each bundle of directions through a hole. Hole selection and a weighted local-lemma estimate rule out triangles, rectangles, and pentagons of lines; the local trades preserve regularity. A uniform bound from Linnik's theorem supplies compatible parameters~\cite[Theorem~1.1]{Xylouris2011}.

Low-degree polynomial products and affine-line restrictions also occur in the coding-theoretic constructions of Golowich--Lin and Dinur--Liu--Zhang~\cite{GolowichLin2024,DinurLiuZhang2025}: degree bounds add under products, while restriction to a line does not increase degree. Here these features produce the line and total pairing identities. We use no coding-theoretic theorem from those works. The only non-elementary existence theorem imported into the direct construction is the uniform least-prime bound; the group embeddings, finite matrix block, interpolation, and probability argument needed here are proved below.

Section~\ref{sec:pairing} first explains the polynomial pairing and the two size bounds. Section~\ref{sec:criterion} formulates the algebraic criterion and proves that its local groups embed. Section~\ref{sec:packing} completes the passage from subgroup averages to scalar elements, including a certificate for the reverse defect. Section~\ref{sec:hypergraph} then supplies the required affine-line system by the full probability and local-trade argument. Section~\ref{sec:prescription} fixes all choices and completes the proof of Theorem~\ref{thm:main}. Section~\ref{sec:cellular} gives the cellular automaton and determines its image by finite local equations.

\paragraph{Conventions.}
All incidence structures in the proof are finite. Their incidence graphs have one vertex for each point and one for each member of the family, with an edge for each membership. The girth is the length of a shortest cycle, and is infinite for a forest. An affine $Q$-line is a translate of a one-dimensional $Q$-subspace. A direction is such a linear subspace, not a chosen nonzero vector. Superscript $\mathsf T$ denotes transpose, without any involution on the coefficient field. All group presentations define abstract discrete groups.
